\section{复现入口与协作清单}
\subsection{源码入口索引}
\begin{longtable}{p{31mm}p{110mm}}
\toprule 内容 & 仓库路径\\\midrule\endhead
最新人形路线 & \code{humanoid/maze_route/marker_route.py}\\
检测与相机 & \code{humanoid/maze_route/markers.py}，\code{robust_camera.py}\\
动作副本 & \code{humanoid/maze_route/actions/}；生成规则见同目录\code{mkcopies.py}\\
通用反馈与道路 & \code{humanoid/scripts/competition.py}，\code{lane_control.py}，\code{lane_trial.py}\\
方式二入口 & \code{turtlebot/follow2/multi_robot_following/launch/vm_follow2.launch}\\
目标与导航控制 & 该包\code{scripts/marker_target.py}、\code{map_follower.py}、\code{velocity_gate.py}\\
轨迹与几何核心 & 该包\code{src/competition_core/planning.py}\\
导航默认参数 & 该包\code{config/move_base.yaml}及\code{costmap_common.yaml}\\\bottomrule
\end{longtable}

\subsection{只查看最新标记检测}
在具备相机、标定文件及显示环境的TonyPi上，从完整仓库运行：
\begin{lstlisting}[language=bash]
cd ~/intelligent_robot_system
python3 humanoid/maze_route/marker_route.py --no-servos
\end{lstlisting}
\code{--no-servos}不访问舵机串口并隐含dry-run；单独\code{--dry-run}仍可能设置头部，二者不能混淆。真实运动入口会关闭厂家上位机并停止相关服务，运行前核对同目录README与设备配置；本报告构建不会运行机器人控制脚本。

\subsection{方式二现场启动索引}
底盘、相机、扫描与地图请按\code{docs/turtlebot_debug_commands.md}逐终端启动，避免重复驱动。完成地图及AMCL初始位姿后，主要跟随入口为：
\begin{lstlisting}[language=bash]
source /opt/ros/kinetic/setup.bash
source ~/catkin_ws/devel/setup.bash
roslaunch multi_robot_following vm_follow2.launch \
  map_file:=$HOME/catkin_ws/maps/competition.yaml \
  marker_id:=1 marker_size_cm:=9.0
\end{lstlisting}
标记尺寸9.0\,cm是该launch默认值，必须按实际ALVAR黑边尺寸替换；与人形航路14\,cm自定义标记无关。正式跟随还需要持续任务状态和各项许可，单独启动此命令不会自动形成全队协同。

\subsection{报告提交前核对}
\begin{enumerate}
\item 各组核对技术章节、实现入口、图片与实验记录；公共格式统一维护。
\item 补齐照片原件、地图/参数版本、每次试验记录与连续运行证据。
\item 将合成测试与实测结果分开，填入样本数、失败次数及实际验收结论。
\item 核对四种方式的独立性、任务状态适配及终点停靠定义。
\item 编译并检查PDF目录、公式、表格和图片；删除或补齐所有待填写项目。
\end{enumerate}
